% LaTeX doc for Overleaf
\documentclass[11pt]{article}
\usepackage[a4paper,margin=1in]{geometry}
\usepackage{amsmath,amssymb,mathtools}
\usepackage{bm}

\newcommand{\E}{\mathbb{E}}
\newcommand{\StopGrad}{\operatorname{StopGrad}}

\begin{document}

\section*{Training procedure: \texttt{feat/common-private-activations-shared-subspace-shuffled-private-loss}}

This note summarizes the objective implemented in the current shuffled-private-loss variant branched from \texttt{feat/common-private-activations-shared-subspace}. The branch keeps the standard DiT flow-matching / velocity-regression loss, and adds:
\begin{itemize}
  \item a \emph{shared-subspace} decomposition computed from one random consecutive window of transformer blocks,
  \item a spatial local-Gram loss on the resulting common component,
  \item a shuffled-pair cosine-squared private loss that compares different layers after shuffling the batch order on one side.
\end{itemize}

\subsection*{1. Repo-specific shapes and notation}
The implementation operates on \emph{patchified} latent tokens.
\begin{itemize}
  \item Clean training inputs are patchified latent tokens
  \[
    \bm{x}_0 \in \mathbb{R}^{B\times N_{\mathrm{tok}}\times D_{\mathrm{in}}},
  \]
  with latent spatial size $32\times 32$, patch size $p=2$, latent channels $C=4$, so
  \[
    N_{\mathrm{tok}} = 16^2 = 256,
    \qquad
    D_{\mathrm{in}} = p^2 C = 16.
  \]
  \item The corresponding unpatchified latent grid is
  \[
    \bm{X}_0 \in \mathbb{R}^{B\times 32\times 32\times 4}.
  \]
  \item The post-block DiT hidden state after block $\ell$ is
  \[
    \bm{A}_\ell \in \mathbb{R}^{B\times N_{\mathrm{tok}}\times D},
    \qquad \ell=1,\dots,L.
  \]
  \item The shared-subspace window size is
  \[
    m = \texttt{layer\_window\_size}.
  \]
  \item The truncated SVD rank is
  \[
    k = \texttt{shared\_subspace\_rank},
    \qquad
    k_{\mathrm{eff}} = \min(k, N_{\mathrm{tok}}, D).
  \]
  \item The private-loss pair budget is
  \[
    K = \texttt{private\_max\_pairs}.
  \]
  If $K\le 0$, the implementation uses all available layer pairs.
  \item The model prediction is denoted by $v_\theta(\cdot)$, conditioned on timestep $\tau\in[0,1]$ and label / conditioning vector $\bm{y}$.
\end{itemize}

\subsection*{2. Main flow-matching diffusion objective}
For each batch, the code samples independent timesteps
\[
  \tau_b \sim \mathcal{U}(0,1),
  \qquad b=1,\dots,B,
\]
and Gaussian noise
\[
  \bm{x}_1 \sim \mathcal{N}(\bm{0}, \bm{I}),
  \qquad
  \bm{x}_1 \in \mathbb{R}^{B\times N_{\mathrm{tok}}\times D_{\mathrm{in}}}.
\]
The forward interpolation uses the repo convention $\tau=0$ for pure noise and $\tau=1$ for clean data:
\[
  \bm{x}_{\tau,b} = (1-\tau_b)\bm{x}_{1,b} + \tau_b \bm{x}_{0,b}.
\]
The regression target is the velocity
\[
  \bm{t}_b = \bm{x}_{0,b} - \bm{x}_{1,b}.
\]
The model predicts
\[
  \hat{\bm{v}}_\theta = v_\theta(\bm{x}_\tau,\bm{\tau},\bm{y}),
\]
and the main loss is
\[
  \ell_{\mathrm{diff}}
  =
  \frac{1}{B N_{\mathrm{tok}} D_{\mathrm{in}}}
  \left\|
    \hat{\bm{v}}_\theta - \bm{t}
  \right\|_2^2.
\]

\paragraph{Implementation detail: internal timestep convention.}
Inside the DiT backbone, the scalar timestep is flipped before the timestep embedder:
\[
  \bar{\tau}_b = 1-\tau_b.
\]
The same flipped convention is reused by the common-spatial projector described later.

\subsection*{3. Random shared window and token-wise normalization}
When auxiliary losses are enabled, the model also returns post-block activations
\[
  \{\bm{A}_\ell\}_{\ell=1}^{L}.
\]
The code uniformly samples a start index $s$ and keeps one consecutive window
\[
  \mathcal{W} = \{s+1, s+2, \dots, s+m\},
\]
where $m$ is clipped to at most $L$.

For each layer in the window, the code applies token-wise LayerNorm across the hidden dimension $D$, without learned affine parameters:
\[
  \tilde{\bm{A}}_\ell[b,n,:]
  =
  \frac{
    \bm{A}_\ell[b,n,:] - \mu_{\ell,b,n}
  }{
    \sqrt{\sigma_{\ell,b,n}^2 + \varepsilon}
  },
\]
where $\mu_{\ell,b,n}$ and $\sigma_{\ell,b,n}^2$ are the mean and variance over the hidden channel index, with $\varepsilon = 10^{-6}$.

\subsection*{4. Shared subspace from truncated SVD}
Within the sampled window, the mean normalized activation is
\[
  \bm{M}
  =
  \frac{1}{m}\sum_{\ell\in\mathcal{W}}\tilde{\bm{A}}_\ell
  \in \mathbb{R}^{B\times N_{\mathrm{tok}}\times D}.
\]
For each batch item $b$, the code computes a truncated right-singular-vector basis from
\[
  \bm{M}_b = \bm{U}_b \bm{\Sigma}_b \bm{V}_b^\top
\]
and keeps the first $k_{\mathrm{eff}}$ right singular vectors:
\[
  \bm{Q}_b
  =
  \left[\bm{v}_{b,1}, \dots, \bm{v}_{b,k_{\mathrm{eff}}}\right]
  \in \mathbb{R}^{D\times k_{\mathrm{eff}}}.
\]

The stop-gradient behavior is controlled by
\[
  \texttt{--shared-subspace-stopgrad}.
\]
When this flag is enabled, the implementation uses
\[
  \bm{M}_b^\dagger = \StopGrad(\bm{M}_b),
  \qquad
  \bm{Q}_b^\dagger = \StopGrad(\bm{Q}_b),
\]
so the SVD path is treated as a detached target extractor. When the flag is disabled, the same formulas below are used but without the detachments.

Equivalently, the shared projector is
\[
  \bm{P}_b = \bm{Q}_b^\sharp (\bm{Q}_b^\sharp)^\top \in \mathbb{R}^{D\times D},
\]
where $\bm{Q}_b^\sharp$ denotes either $\bm{Q}_b^\dagger$ or $\bm{Q}_b$ depending on the CLI toggle.
The implementation does not explicitly materialize $\bm{P}_b$; it applies the projector as
\[
  \tilde{\bm{A}}_\ell \bm{P}
  =
  (\tilde{\bm{A}}_\ell \bm{Q}^\sharp)(\bm{Q}^\sharp)^\top.
\]

\subsection*{5. Common decomposition from the shared window}
For each layer $\ell\in\mathcal{W}$ inside the sampled consecutive window,
\[
  \bm{A}^{\mathrm{com}}_\ell[b,:,:]
  =
  \tilde{\bm{A}}_\ell[b,:,:]\bm{P}_b.
\]
The common activation used by the auxiliary branches is the mean over the same window:
\[
  \bm{A}^{\mathrm{com}}_{\mathrm{mean}}
  =
  \frac{1}{m}\sum_{\ell\in\mathcal{W}} \bm{A}^{\mathrm{com}}_\ell
  \in
  \mathbb{R}^{B\times N_{\mathrm{tok}}\times D}.
\]
This tensor is the shared reference used by both:
\begin{itemize}
  \item the spatial branch, which projects it into token-space features,
  \item the private branch, which subtracts it from other randomly chosen layers.
\end{itemize}

\subsection*{6. Spatial branch}
\subsubsection*{6.1 Common-to-token projector}
The spatial loss compares the common representation to a coarse latent target in the same token-feature space. The branch maps
\[
  \bm{A}^{\mathrm{com}}_{\mathrm{mean}}
  \in \mathbb{R}^{B\times 256\times D}
\]
through
\[
  \Pi_{\mathrm{sp}}
  =
  \Pi_{\tau}
  \circ
  \Pi_{\mathrm{base}},
\]
where:
\begin{itemize}
  \item $\Pi_{\mathrm{base}}$ is controlled by \texttt{common\_spatial\_projector}; it is either identity or a lightweight CNN on the $16\times 16$ token grid followed by a linear map back to hidden size $D$,
  \item $\Pi_{\tau}$ applies a token-wise LayerNorm and then timestep-conditioned modulation
  \[
    \Pi_{\tau}(\bm{H},\tau_b)[n,:]
    =
    \operatorname{LN}(\bm{H}[n,:]) \odot (1+\bm{s}_b) + \bm{h}_b,
  \]
  where $\bm{s}_b$ and $\bm{h}_b$ are produced from the flipped timestep embedding $e_t(1-\tau_b)$.
\end{itemize}
In the code, the shift head for $\bm{h}_b$ is zero-initialized, while the scale head uses the default Dense initialization.

The projected common tensor is
\[
  \bm{Z}
  =
  \Pi_{\mathrm{sp}}\!\left(\bm{A}^{\mathrm{com}}_{\mathrm{mean}}, \bm{\tau}\right)
  \in \mathbb{R}^{B\times 256\times D}
  \cong
  \mathbb{R}^{B\times 16\times 16\times D}.
\]

\subsubsection*{6.2 Coarse timestep-interpolated target}
The clean latent tokens $\bm{x}_0$ are first unpatchified to
\[
  \bm{X}_0 \in \mathbb{R}^{B\times 32\times 32\times 4}.
\]
The code average-pools $\bm{X}_0$ down to $4\times 4$:
\[
  \bm{X}_{\mathrm{coarse}}
  =
  \mathrm{AvgPool}_{8\times 8,\;\mathrm{stride}=8}(\bm{X}_0),
\]
then upsamples bilinearly back to $32\times 32$:
\[
  \bm{X}_{\mathrm{low}}
  =
  \mathrm{BilinearResize}(\bm{X}_{\mathrm{coarse}}, 32, 32).
\]
The timestep-interpolated target is
\[
  \bm{X}_{\mathrm{tgt},b}
  =
  (1-\tau_b)\bm{X}_{\mathrm{low},b}
  +
  \tau_b \bm{X}_{0,b}.
\]
This target is patchified back into
\[
  \bm{X}^{\mathrm{patch}}_{\mathrm{tgt}}
  \in
  \mathbb{R}^{B\times 256\times 16},
\]
then mapped to hidden dimension only through the DiT patch-embedding layer:
\[
  \widehat{\bm{X}}_{\mathrm{tgt}}
  =
  E_{\mathrm{patch}}\!\left(\bm{X}^{\mathrm{patch}}_{\mathrm{tgt}}\right)
  \in
  \mathbb{R}^{B\times 256\times D}.
\]
Unlike the main DiT forward pass, this target branch does \emph{not} add positional embeddings. Before entering the Gram loss, it applies the same token-wise LayerNorm as in Section 3:
\[
  \widetilde{\bm{X}}_{\mathrm{tgt}}
  =
  \operatorname{LN}_{\mathrm{tok}}\!\left(\widehat{\bm{X}}_{\mathrm{tgt}}\right)
  \in
  \mathbb{R}^{B\times 256\times D}
  \cong
  \mathbb{R}^{B\times 16\times 16\times D}.
\]
Therefore both spatial branches live in the same token-grid feature space:
\[
  \bm{Z},\ \widetilde{\bm{X}}_{\mathrm{tgt}}
  \in
  \mathbb{R}^{B\times 16\times 16\times D}.
\]

\subsection*{7. Local Gram spatial loss}
Let
\[
  s = \texttt{spatial\_window\_size},
  \qquad
  \Delta = \texttt{spatial\_window\_stride}.
\]
Both $\bm{Z}$ and $\widetilde{\bm{X}}_{\mathrm{tgt}}$ are decomposed into all valid sliding windows of size $s\times s$ over the $16\times 16$ token grid. One window contains
\[
  M = s^2
\]
token positions. The number of valid window locations per axis is
\[
  n_{\Omega}
  =
  \left\lfloor \frac{16-s}{\Delta}\right\rfloor + 1,
\]
so the total number of windows is
\[
  |\Omega| = n_{\Omega}^2.
\]
After extraction, both tensors have shape
\[
  \bm{Z}_{\mathrm{win}},
  \widetilde{\bm{X}}^{\mathrm{win}}_{\mathrm{tgt}}
  \in
  \mathbb{R}^{B\times |\Omega|\times M\times D}.
\]

For each batch item $b$, window $w$, and token-position indices $u,v\in\{1,\dots,M\}$, the local Gram entries are
\[
  G^{\mathrm{feat}}_{b,w,u,v}
  =
  \frac{1}{D}\sum_{c=1}^{D}
  Z_{b,w,u,c}\,Z_{b,w,v,c},
\]
\[
  G^{\mathrm{tgt}}_{b,w,u,v}
  =
  \frac{1}{D}\sum_{c=1}^{D}
  \widetilde{X}^{\mathrm{tgt}}_{b,w,u,c}\,\widetilde{X}^{\mathrm{tgt}}_{b,w,v,c}.
\]
Equivalently, for fixed $(b,w)$,
\[
  \bm{G}^{\mathrm{feat}}_{b,w}
  =
  \frac{1}{D}\bm{Z}_{b,w}\bm{Z}_{b,w}^{\top},
  \qquad
  \bm{G}^{\mathrm{tgt}}_{b,w}
  =
  \frac{1}{D}\widetilde{\bm{X}}^{\mathrm{tgt}}_{b,w}
  \left(\widetilde{\bm{X}}^{\mathrm{tgt}}_{b,w}\right)^{\top}.
\]
The target branch is detached. The spatial loss is
\[
  \ell_{\mathrm{spatial}}
  =
  \frac{1}{B |\Omega| M^2}
  \sum_{b=1}^{B}
  \sum_{w\in\Omega}
  \sum_{u=1}^{M}
  \sum_{v=1}^{M}
  \left|
    G^{\mathrm{feat}}_{b,w,u,v}
    -
    \StopGrad\!\left(G^{\mathrm{tgt}}_{b,w,u,v}\right)
  \right|.
\]
The implementation computes this exactly but uses only upper-triangular Gram entries together with diagonal/off-diagonal weights to avoid materializing the redundant lower triangle.

\subsection*{8. Shuffled-pair private loss}
The private loss does \emph{not} reuse all pairwise comparisons inside the consecutive shared window. Instead, after constructing $\bm{A}^{\mathrm{com}}_{\mathrm{mean}}$, the code samples random layer pairs from the \emph{full} depth:
\[
  \mathcal{P}_{\mathrm{all}}
  =
  \{(\ell_i,\ell_j): 1\le \ell_i < \ell_j \le L\}.
\]
Let
\[
  K_{\mathrm{eff}}
  =
  \begin{cases}
    |\mathcal{P}_{\mathrm{all}}|, & K \le 0,\\
    \min(K, |\mathcal{P}_{\mathrm{all}}|), & K > 0.
  \end{cases}
\]
The code samples
\[
  \mathcal{P} = \{(\ell_i,\ell_j)\}_{p=1}^{K_{\mathrm{eff}}}
  \subseteq
  \mathcal{P}_{\mathrm{all}}
\]
uniformly without replacement.

For each sampled pair $(\ell_i,\ell_j)\in\mathcal{P}$, it computes token-wise normalized activations
\[
  \tilde{\bm{A}}_{\ell_i},\ \tilde{\bm{A}}_{\ell_j}
  \in \mathbb{R}^{B\times N_{\mathrm{tok}}\times D},
\]
then subtracts the same shared common tensor:
\[
  \bm{R}^{\mathrm{L}}_{p}
  =
  \tilde{\bm{A}}_{\ell_i} - \bm{A}^{\mathrm{com}}_{\mathrm{mean}},
  \qquad
  \bm{R}^{\mathrm{R}}_{p}
  =
  \tilde{\bm{A}}_{\ell_j} - \bm{A}^{\mathrm{com}}_{\mathrm{mean}}.
\]
Thus the private branch measures what remains in randomly chosen layers after removing the common component distilled from the consecutive window.

Each side is flattened over both token and hidden dimensions:
\[
  \bm{f}^{\mathrm{L}}_{p}[b,:]
  =
  \operatorname{vec}\!\bigl(\bm{R}^{\mathrm{L}}_{p}[b,:,:]\bigr)
  \in \mathbb{R}^{N_{\mathrm{tok}}D},
  \qquad
  \bm{f}^{\mathrm{R}}_{p}[b,:]
  =
  \operatorname{vec}\!\bigl(\bm{R}^{\mathrm{R}}_{p}[b,:,:]\bigr)
  \in \mathbb{R}^{N_{\mathrm{tok}}D},
\]
and then $\ell_2$-normalized:
\[
  \hat{\bm{f}}^{\mathrm{L}}_{p}[b,:]
  =
  \frac{\bm{f}^{\mathrm{L}}_{p}[b,:]}
       {\|\bm{f}^{\mathrm{L}}_{p}[b,:]\|_2+\varepsilon},
  \qquad
  \hat{\bm{f}}^{\mathrm{R}}_{p}[b,:]
  =
  \frac{\bm{f}^{\mathrm{R}}_{p}[b,:]}
       {\|\bm{f}^{\mathrm{R}}_{p}[b,:]\|_2+\varepsilon}.
\]

For each sampled pair $p$, the code draws a non-zero cyclic batch shift
\[
  \delta_p \sim \operatorname{Unif}\{1,\dots,B-1\},
\]
and defines
\[
  \pi_p(b) = ((b-1)+\delta_p)\bmod B + 1.
\]
Therefore
\[
  \pi_p(b)\neq b
\]
for every batch index $b$, so the right-hand side is always matched against a \emph{different} batch element. The cosine used by the loss is
\[
  c_p(b)
  =
  \hat{\bm{f}}^{\mathrm{L}}_{p}[b,:]^\top
  \hat{\bm{f}}^{\mathrm{R}}_{p}[\pi_p(b),:].
\]
The private loss is the batch-and-pair average of squared cosines:
\[
  \ell_{\mathrm{private}}
  =
  \frac{1}{B K_{\mathrm{eff}}}
  \sum_{p=1}^{K_{\mathrm{eff}}}
  \sum_{b=1}^{B}
  c_p(b)^2.
\]
Hence the objective explicitly pushes a private summary from one layer to differ from a shuffled counterpart that comes from:
\begin{itemize}
  \item another layer,
  \item another batch element,
  \item and, in ordinary training, therefore typically another image and another sampled timestep.
\end{itemize}
If $L<2$, no layer pairs exist and the private loss is zero. If $B<2$, the code also returns zero private loss because no non-trivial batch shuffle is possible.

\subsection*{9. Total loss and schedules}
The total objective is
\[
  \mathcal{L}
  =
  \ell_{\mathrm{diff}}
  +
  \lambda_{\mathrm{spatial}}^{\mathrm{eff}} \ell_{\mathrm{spatial}}
  +
  \lambda_{\mathrm{private}}^{\mathrm{eff}} \ell_{\mathrm{private}}.
\]
In the CLI, the private coefficient uses \texttt{--lambda-private} and also accepts the alias \texttt{--lambda-div}.

\subsubsection*{9.1 Private-loss warmup}
Let $t$ be the current training step, $t_0=\texttt{private\_start\_step}$, and $T=\texttt{private\_warmup\_iters}$. Then
\[
  \lambda_{\mathrm{private}}^{\mathrm{eff}}(t)
  =
  \lambda_{\mathrm{private}} \cdot \gamma_{\mathrm{private}}(t),
\]
with
\[
  \gamma_{\mathrm{private}}(t)
  =
  \begin{cases}
    0, & t < t_0,\\[4pt]
    1, & t\ge t_0 \text{ and } T\le 0,\\[4pt]
    \mathrm{clip}\!\left(\dfrac{t-t_0+1}{T}, 0, 1\right),
    & t\ge t_0 \text{ and } T>0.
  \end{cases}
\]

\subsubsection*{9.2 Spatial-loss decay}
Let $t_{\mathrm{stop}}=\texttt{spatial\_stop\_step}$ and $T_{\mathrm{stop}}=\texttt{spatial\_stop\_warmup\_iters}$. Then
\[
  \lambda_{\mathrm{spatial}}^{\mathrm{eff}}(t)
  =
  \lambda_{\mathrm{spatial}} \cdot \gamma_{\mathrm{spatial}}(t),
\]
where
\[
  \gamma_{\mathrm{spatial}}(t)
  =
  \begin{cases}
    1, & t_{\mathrm{stop}} < 0,\\[4pt]
    1, & t < t_{\mathrm{stop}},\\[4pt]
    0, & t\ge t_{\mathrm{stop}} \text{ and } T_{\mathrm{stop}}\le 0,\\[4pt]
    1-\mathrm{clip}\!\left(\dfrac{t-t_{\mathrm{stop}}+1}{T_{\mathrm{stop}}}, 0, 1\right),
    & t\ge t_{\mathrm{stop}} \text{ and } T_{\mathrm{stop}}>0.
  \end{cases}
\]

\subsection*{10. High-level algorithm per training step}
\begin{enumerate}
  \item Sample per-example timesteps $\tau_b\sim\mathcal{U}(0,1)$ and Gaussian noise $\bm{x}_1\sim\mathcal{N}(\bm{0},\bm{I})$.
  \item Build
  \[
    \bm{x}_\tau = (1-\bm{\tau})\bm{x}_1 + \bm{\tau}\bm{x}_0,
    \qquad
    \bm{t} = \bm{x}_0 - \bm{x}_1.
  \]
  \item Run the DiT on $(\bm{x}_\tau,\bm{\tau},\bm{y})$ to get $\hat{\bm{v}}_\theta$ and block activations $\{\bm{A}_\ell\}_{\ell=1}^{L}$.
  \item Uniformly sample one consecutive window $\mathcal{W}$ of size $m$, token-normalize it, average it, and compute a rank-$k$ shared subspace via truncated SVD. Depending on \texttt{--shared-subspace-stopgrad}, this SVD path is either detached or fully differentiable.
  \item Project the normalized window activations onto the shared subspace and average them to obtain $\bm{A}^{\mathrm{com}}_{\mathrm{mean}}$.
  \item Project $\bm{A}^{\mathrm{com}}_{\mathrm{mean}}$ through the common-spatial head, build the coarse timestep-interpolated target, embed it with the patch embedder, token-normalize it, and compute the local Gram loss in token-feature space.
  \item Independently sample random layer pairs from the full depth, token-normalize both sides, subtract $\bm{A}^{\mathrm{com}}_{\mathrm{mean}}$, flatten each residual over $N_{\mathrm{tok}}D$, batch-shuffle one side with a non-zero cyclic shift, and compute the cosine-squared private loss.
  \item Form
  \[
    \mathcal{L}
    =
    \ell_{\mathrm{diff}}
    +
    \lambda_{\mathrm{spatial}}^{\mathrm{eff}} \ell_{\mathrm{spatial}}
    +
    \lambda_{\mathrm{private}}^{\mathrm{eff}} \ell_{\mathrm{private}},
  \]
  then backpropagate through the online model parameters.
\end{enumerate}

\end{document}
